\documentclass[11pt]{article}

\usepackage[a4paper,margin=1in]{geometry}
\usepackage{hyperref}
\usepackage{enumitem}
\usepackage{titlesec}

\titleformat{\section}{\large\bfseries}{\thesection}{1em}{}
\setlist[itemize]{noitemsep, topsep=4pt}
\setlist[enumerate]{noitemsep, topsep=4pt}

\title{KCPC-CS111 \\ Supplementary Reading \\ Week 2: Techniques for Investigation}
\author{King's Competitive Programming Club}
\date{}

\begin{document}

\maketitle

\noindent
This document contains recommended reading for Workshop 2: Techniques for Investigation. These texts deepen your mathematical problem solving toolkit, spanning heuristics, induction, spatial visualisation, invariance, symmetry, the extreme principle, and bitwise competitive programming.

\vspace{1em}
\hrule
\vspace{1em}

\section*{I. Core Problem-Solving Heuristics}

\subsection*{The Art and Craft of Problem Solving (Paul Zeitz)}
John Wiley \& Sons, 2nd Edition, 2006.

\begin{itemize}
    \item \textbf{Chapter 2: Strategies for Investigation} Psychological strategies, wishful thinking, ``make it easier'', drawing diagrams, and identifying patterns from small examples.
    \item \textbf{Chapter 3: Tactics for Solving Problems} Symmetry (geometric and algebraic), parity, and the extreme principle.
    \item \textbf{Chapter 7: Modular Arithmetic \& Functional Equations} Techniques for substituting zero or special values into multi-variable equations.
\end{itemize}

\subsection*{How to Solve It: A New Aspect of Mathematical Method (George P{\'o}lya)}
Princeton University Press, 1945 / 2004.

\begin{itemize}
    \item \textbf{Part I: In the Classroom} Understanding the problem, devising a plan, carrying out the plan, and looking back.
    \item \textbf{Heuristic Dictionary} Analogy, auxiliary problems, decomposing and recombining, inductive vs. deductive discovery, and specialization.
\end{itemize}

\vspace{1em}
\hrule
\vspace{1em}

\section*{II. The Extreme Principle and Invariants}

\subsection*{Problem-Solving Strategies (Arthur Engel)}
Springer, 1998.

\begin{itemize}
    \item \textbf{Chapter 2: Coloring Proofs} Grid tilings, dominoes, trominoes, and invariant checkerboard colorings.
    \item \textbf{Chapter 3: The Extremal Principle} Choosing minimal and maximal elements in finite graphs, discrete geometries, and number systems to produce immediate contradictions.
\end{itemize}

\vspace{1em}
\hrule
\vspace{1em}

\section*{III. Bitwise Algorithms and Fibonacci Bases}

\subsection*{CP-Algorithms: Bit Manipulation and Bitmasking}
\url{https://cp-algorithms.com/algebra/all-submasks.html}

\begin{itemize}
    \item Bitwise operations ($\&, |, \oplus$), population count (\texttt{\_\_builtin\_popcountll}), and iterating over submasks in $O(3^N)$.
    \item Structural explanation of $f(2n) = f(n)$ and $f(2n+1) = f(n)+1$ as bit-shift operations.
\end{itemize}

\subsection*{Competitive Programmer's Handbook (Antti Laaksonen)}
\url{https://cses.fi/book/book.pdf}

\begin{itemize}
    \item \textbf{Chapter 10: Bit Manipulation} Bit representations, bitwise operations, and dynamic programming with bitmasks.
    \item \textbf{Chapter 21: Number Theory} Fibonacci properties, Zeckendorf representation, and fast matrix exponentiation.
\end{itemize}

\end{document}
